\documentclass[10pt,a4paper]{amsart}
\usepackage{amsmath,amssymb,amsthm,graphicx}
\usepackage{verbatim}
\newenvironment{reference}{\comment}{\endcomment}
%\newenvironment{reference}{}{}
\newenvironment{slogan}{\comment}{\endcomment}
\newenvironment{history}{\comment}{\endcomment}

% For commutative diagrams we use Xy-pic
\usepackage[all]{xy}

% We use 2cell for 2-commutative diagrams.
\xyoption{2cell}
% Theorem environments.
%
\theoremstyle{plain}
\newtheorem{theorem}[subsection]{Theorem}
\newtheorem{proposition}[subsection]{Proposition}
\newtheorem{lemma}[subsection]{Lemma}

\theoremstyle{definition}
\newtheorem{definition}[subsection]{Definition}
\newtheorem{example}[subsection]{Example}
\newtheorem{exercise}[subsection]{Exercise}
\newtheorem{situation}[subsection]{Situation}

\theoremstyle{remark}
\newtheorem{remark}[subsection]{Remark}
\newtheorem{remarks}[subsection]{Remarks}

\numberwithin{equation}{subsection}

% Macros
%
\def\lim{\mathop{\mathrm{lim}}\nolimits}
\def\colim{\mathop{\mathrm{colim}}\nolimits}
\def\Spec{\mathop{\mathrm{Spec}}}
\def\Hom{\mathop{\mathrm{Hom}}\nolimits}
\def\Ext{\mathop{\mathrm{Ext}}\nolimits}
\def\SheafHom{\mathop{\mathcal{H}\!\mathit{om}}\nolimits}
\def\SheafExt{\mathop{\mathcal{E}\!\mathit{xt}}\nolimits}
\def\Sch{\mathit{Sch}}
\def\Mor{\mathop{\mathrm{Mor}}\nolimits}
\def\Ob{\mathop{\mathrm{Ob}}\nolimits}
\def\Sh{\mathop{\mathit{Sh}}\nolimits}
\def\NL{\mathop{N\!L}\nolimits}
\def\CH{\mathop{\mathrm{CH}}\nolimits}
\def\proetale{{pro\text{-}\acute{e}tale}}
\def\etale{{\acute{e}tale}}
\def\QCoh{\mathit{QCoh}}
\def\Ker{\mathop{\mathrm{Ker}}}
\def\Im{\mathop{\mathrm{Im}}}
\def\Coker{\mathop{\mathrm{Coker}}}
\def\Coim{\mathop{\mathrm{Coim}}}

% Boxtimes
%
\DeclareMathSymbol{\boxtimes}{\mathbin}{AMSa}{"02}

%
% Macros for moduli stacks/spaces
%
\def\QCohstack{\mathcal{QC}\!\mathit{oh}}
\def\Cohstack{\mathcal{C}\!\mathit{oh}}
\def\Spacesstack{\mathcal{S}\!\mathit{paces}}
\def\Quotfunctor{\mathrm{Quot}}
\def\Hilbfunctor{\mathrm{Hilb}}
\def\Curvesstack{\mathcal{C}\!\mathit{urves}}
\def\Polarizedstack{\mathcal{P}\!\mathit{olarized}}
\def\Complexesstack{\mathcal{C}\!\mathit{omplexes}}
% \Pic is the operator that assigns to X its picard group, usage \Pic(X)
% \Picardstack_{X/B} denotes the Picard stack of X over B
% \Picardfunctor_{X/B} denotes the Picard functor of X over B
\def\Pic{\mathop{\mathrm{Pic}}\nolimits}
\def\Picardstack{\mathcal{P}\!\mathit{ic}}
\def\Picardfunctor{\mathrm{Pic}}
\def\Deformationcategory{\mathcal{D}\!\mathit{ef}}
\usepackage{cleveref}
\crefname{originallemma}{original lemma}{original lemmas}
\Crefname{originallemma}{Original lemma}{Original lemmas}
\makeatletter
\newlabel{fields-lemma-transcendence-degree}{{030F}{source}{}{sourcepointer-1}{}}
\newlabel{fields-lemma-transcendence-degree@cref}{{[originallemma][0][]030F}{[1][0][]source}{}{}{}}
\makeatother
\title{A separating transcendence base exists among the generators of the extension}
\author{The Stacks Project Authors}
\date{stacks-project:algebra.tex:local-snapshot-20261012; GNU FDL 1.2 or later}
\begin{document}
\maketitle
\noindent\textit{Editorial note.} All delivered mathematical text is the author's own, copied from the pinned Stacks Project chapter file algebra.tex; no statement or proof step was rewritten, reordered or removed.\par\medskip
\noindent\textit{Editorial note.} Preamble provenance, as required for this source: the frozen snapshot stages the chapter's own preamble file beside it (preamble.tex, 217 lines, 8116 bytes, the same shared Stacks-wide preamble the published chow and etale-cohomology cards use); the first line of the pinned chapter file is \textbackslash{}input\{preamble\}, so no preamble text lives inside the chapter itself. The card takes that file's lines 8-18 (the verbatim package and the three \textbackslash{}comment-based environments reference/slogan/history at lines 9-12, section 13.2, together with the Xy-pic setup) and its lines 154 to the runtime end of the file (the native theorem declarations and mathematical macros). Not one character was changed, no package was added, no macro was defined, and the excerpt serves only as preamble, never as a span coordinate.\par\medskip
\noindent\textit{Editorial note.} The same-chapter pointers are located at their \textbackslash{}label lines inside the pinned chapter file; the cross-chapter pointers are located in the project's own tag list staged in the snapshot (sources/stacks-chow/tags.txt), at the verbatim lines carrying their canonical tags.\par\medskip
\noindent\textit{Editorial note.} The counted proof refers to same-chapter statements by label; they are recorded as navigation stubs pointed at their label lines.\par\medskip
\noindent\textit{Editorial note.} The counted spans contain no cite, so the card carries no bibliography.\par\medskip
\noindent\textit{Allowed standard prerequisites (editorial).}\begin{itemize}
\item Minimal polynomial of a generator over the rational function field. Used at the start of the counted proof (lines 10528-10532). Produces the irreducible relation whose coefficients drive the argument.
\item Separability tested by non-containment in the p-th power subring. Used at the end of the counted proof (lines 10564-10566). Converts the non-containment of the polynomial into separability of the extension.
\end{itemize}\medskip
\ifdefined\hypertarget\hypertarget{sourcepointer-1}{}\fi
\noindent\textit{Source reference.} Original lemma 030F: sources/stacks-chow/tags.txt, source lines 3946--3946. Navigation pointer for \textbackslash{}ref\{fields-lemma-transcendence-degree\} used inside the counted span (line 10555); the argument that consumes it is retained in the span.\par\medskip
\begin{lemma}
\label{lemma-mini-separability}
Let $k$ be a field of characteristic $p > 1$. Let $K/k$
be a field extension generated by $x_1, \ldots, x_{n + 1} \in K$ such that
\begin{enumerate}
\item $\{x_1, \ldots, x_n\}$ is a transcendence base of $K/k$,
\item for every $k$-linearly independent subset
$\{a_1, \ldots, a_m\}$ of $K$ the set
$\{a^p_1, \ldots, a_m^p\}$ is $k$-linearly independent.
\end{enumerate}
Then there is $1 \leq j \leq n+1$ such that
$\{ x_1, \ldots, \widehat{x}_j, \ldots, x_{n+1}\}$
is a separating transcendence base for $K / k$.
\end{lemma}

\begin{proof}
By assumption $x_{n + 1}$ is algebraic over $k(x_1, \ldots, x_n)$
so there exists a non-zero polynomial $F \in k[X_1, \ldots, X_{n + 1}]$
such that $F(x_1, \ldots, x_{n+1}) = 0$. Choose $F$ of minimal total degree.
Then $F$ is irreducible, because at least one irreducible factor
must also have the same property.

\medskip\noindent
We claim that, for some $i$, not all powers of $X_i$ appearing in $F$
are multiples of $p$. Suppose for a contradiction that all the exponents
appearing in $F$ were multiples of $p$, then the set
$$
\{x_1^{\alpha_1} \ldots x^{\alpha_{n+1}}_{n+1} \mid \lambda_\alpha \neq 0\}
\subset
K
$$
is $k$-linearly dependent where $\lambda_\alpha$ are the coefficients of $F$.
By assumption (2) we conclude the set
$$
\{x_1^{\alpha_1 / p} \ldots x^{\alpha_{n+1} / p}_{n+1}
\mid \lambda_\alpha \neq 0 \}
$$
is also $k$-linearly dependent, contradicting minimality of $\deg(F)$.

\medskip\noindent
Choose $i$ for which a non-$p$th power of $X_i$ appears in $F$. Then we
see that $x_i$ is algebraic over
$L = k(x_1, \ldots, x_{i - 1}, x_{i + 1}, \ldots, x_{n+1})$.
By Fields, Lemma \ref{fields-lemma-transcendence-degree}
we see that $x_1, \ldots, x_{i - 1}, x_{i + 1}, \ldots, x_{n+1}$
is a transcendence base of $K/k$. Thus $L$ is the fraction field
of the polynomial ring over $k$ in
$x_1, \ldots, x_{i - 1}, x_{i + 1}, \ldots, x_{n + 1}$.
By Gauss' Lemma we conclude that
$$
P(T) =
F(x_1, \ldots, x_{i - 1}, T, x_{i + 1}, \ldots, x_{n + 1}) \in L[T]
$$
is irreducible. By construction $P(T)$ is not contained in $L[T^p]$.
Hence $K/L$ is separable as required.
\end{proof}

\end{document}
